\section{A Saturação da Engenharia de Características Clássica ($H_3$)}
\label{sec:saturacao_features}

A refutação da Hipótese $H_3$, atestada pela degradação massiva de performance durante a Fase 3, fornece uma contribuição teórica relevante sobre a evolução e os limites dos paradigmas na visão computacional. Historicamente, a engenharia de características manuais (\textit{feature engineering}), como a aplicação de bancos de filtros de Gabor, representou o estado da arte para a extração de assinaturas de texturas finas. Contudo, os resultados empíricos demonstraram que, quando acoplada a arquiteturas profundas modernas, a injeção estática de conhecimento morfológico não apenas se revela redundante, como atua de forma severamente deletéria.

O colapso preditivo observado ao expandir o tensor de entrada de 3 canais brutos (RGB) para 15 canais filtrados pode ser explicado pelos fenômenos do ruído paramétrico e da saturação representacional. O paradigma do aprendizado de ponta a ponta pressupõe que os pesos sinápticos são progressivamente ajustados para descobrir as representações latentes mais eficientes para uma determinada tarefa. Na literatura consolidada, sabe-se que as camadas convolucionais iniciais de redes profundas convergem naturalmente para detectores de bordas, linhas e estrias que são matematicamente análogos aos filtros de Gabor, porém otimizados de forma dinâmica para a minimização da função de perda específica do problema.

Ao forçar os classificadores a operar sobre um conjunto de dados de representações espectrais pré-computadas e estáticas, o desenho experimental limitou o espaço de busca na otimização estocástica. A multiplicidade de canais altamente correlacionados induziu uma espécie de interferência destrutiva, diluindo o sinal físico original e engessando o fluxo de gradientes. Em suma, as redes despenderam sua capacidade paramétrica tentando decodificar um ruído estrutural induzido, em vez de focar na fenotipagem do algodão. É importante notar que a normalização conjunta dos 15 canais pode ter induzido um achatamento na variância dos pixels originais, criando um gargalo representacional que as camadas profundas não conseguiram reverter, resultando no colapso do gradiente observado durante a Fase 3.

Este achado valida a premissa fundamental do \textit{Deep Learning} aplicado à agroindústria: a arquitetura convolucional possui uma capacidade de abstração intrínseca superior à modelagem matemática geométrica clássica. A constatação consagra o entendimento de que, uma vez garantida a estabilidade fotométrica na captura geométrica do problema (conforme provado em $H_2$), a pureza do sinal multidimensional bruto (RGB) é a condição que maximiza a assertividade da inteligência artificial.